\section{Вероятности экстремальных значений}

Пусть $X_1, X_2, \dots, X_n$ --- независимые одинаково распределенные случайные величины с функцией распределения $F_\xi(x) = P(\xi < x)$.

\begin{definition}[Максимум и Минимум]
\leavevmode 
\begin{itemize}
    \item $M_n = \max(X_1, \dots, X_n)$
    
    Функция распределения максимума:
    \[ F_{M_n}(x) = P(M_n < x) = P(X_1 < x, \dots, X_n < x) = (F_\xi(x))^n \]
    
    \item $\tilde{M}_n = \min(X_1, \dots, X_n)$
    
    Функция распределения минимума (через дополнение):
    \[ 1 - F_{\tilde{M}_n}(x) = P(\tilde{M}_n \ge x) = P(X_1 \ge x, \dots, X_n \ge x) = (1 - F_\xi(x))^n \]
\end{itemize}
\end{definition}

\begin{example}[Равномерное распределение]
Пусть $X_i \sim U[a, b]$.
$$ F_\xi(x) = \begin{cases} 0, & x < a \\ \frac{x-a}{b-a}, & x \in [a, b] \\ 1, & x > b \end{cases} $$
Тогда для максимума $M_n$:
$$ F_{M_n}(x) = \left( \frac{x-a}{b-a} \right)^n, \quad x \in [a, b] $$
Плотность распределения $f_{M_n}(x) = F'_{M_n}(x) = \frac{n(x-a)^{n-1}}{(b-a)^n}$.

Математическое ожидание максимума:
$$ E[M_n] = \int_a^b x \cdot \frac{n(x-a)^{n-1}}{(b-a)^n} dx $$
Замена $y = x-a$, $dy=dx$, пределы $0 \dots b-a$:
$$ E[M_n] = \frac{n}{(b-a)^n} \int_0^{b-a} (y+a)y^{n-1} dy = \frac{n}{(b-a)^n} \left[ \frac{y^{n+1}}{n+1} + a\frac{y^n}{n} \right]_0^{b-a} $$
$$ = \frac{n}{(b-a)^n} \left( \frac{(b-a)^{n+1}}{n+1} + a\frac{(b-a)^n}{n} \right) = \frac{n(b-a)}{n+1} + a = \frac{nb + a}{n+1} $$
\end{example}


\subsection{Предельные теоремы для экстремальных значений}

\begin{theorem}[Центральная предельная теорема (напоминание)]
Пусть $\xi_1, \dots, \xi_n$ --- независимые одинаково распределенные случайные величины с $E\xi = a$, $D\xi = \sigma^2$.
Тогда:
$$ P\left(\frac{S_n - na}{\sqrt{n}\sigma} \le x\right) \to \Phi(x) $$
где $S_n = \xi_1 + \dots + \xi_n$.
\end{theorem}

\begin{theorem}[Предельное распределение максимума]
\leavevmode
Пусть $M_n = \max(\xi_1, \dots, \xi_n)$, где $\xi_i$ независимы и одинаково распределены с функцией распределения $F(x)$.
Если существует последовательность $u_n(x)$ такая, что:
$$ n(1 - F(u_n)) \xrightarrow{n \to \infty} x, \quad 0 < x < \infty $$
то:
$$ P(M_n \le u_n(x)) \xrightarrow{n \to \infty} e^{-x} $$
\end{theorem}

\begin{proof}
Из условия: $n(1 - F(u_n)) = x + o(1)$, откуда:
$$ 1 - F(u_n) = \frac{x}{n} + o\left(\frac{1}{n}\right) $$
$$ F(u_n) = 1 - \frac{x}{n} + o\left(\frac{1}{n}\right) $$

Тогда:
$$ P(M_n \le u_n(x)) = (F(u_n(x)))^n = \left(1 - \frac{x}{n} + o\left(\frac{1}{n}\right)\right)^n $$
$$ = \exp\left[n \cdot \ln\left(1 - \frac{x}{n} + o\left(\frac{1}{n}\right)\right)\right] = e^{-x} \cdot e^{o(1)} = e^{-x}(1 + o(1)) $$
\end{proof}

\begin{example}[Экспоненциальное распределение]
Пусть $F(x) = 1 - e^{-\lambda x}$.
Найдем нормирующую последовательность $u_n(\tau)$ из условия:
$$ n(1 - F(u_n(\tau))) \to \tau $$
$$ n \cdot e^{-\lambda u_n(\tau)} = \tau(1 + o(1)) $$
$$ e^{-\lambda u_n(\tau)} = \frac{\tau}{n}(1 + o(1)) $$
$$ -\lambda u_n(\tau) = \ln\tau - \ln n + \ln(1 + o(1)) $$
$$ u_n(\tau) = \frac{1}{\lambda}(\ln n - \ln\tau) $$

Тогда:
$$ P\left(M_n \le \frac{1}{\lambda}(\ln n - \ln\tau)\right) = e^{-\tau} + o(1) $$
$$ P\left(M_n - \frac{\ln n}{\lambda} \le -\frac{\ln\tau}{\lambda}\right) = e^{-\tau} + o(1) $$

Замена $y = -\frac{\ln\tau}{\lambda}$, тогда $\tau = e^{-\lambda y}$:
$$ P\left(\lambda M_n - \ln n \le y\right) = e^{-e^{-y}} + o(1) $$
Это распределение Гумбеля.
\end{example}

\begin{example}[Равномерное распределение на {[0, 1]}]
$$ F(x) = \begin{cases} 0, & x \le 0 \\ x, & x \in (0, 1) \\ 1, & x \ge 1 \end{cases} $$

Из условия $n(1 - F(u_n(\tau))) \to \tau$:
$$ n(1 - u_n(\tau)) = \tau $$
$$ u_n(\tau) = 1 - \frac{\tau}{n} $$

Тогда:
$$ P\left(M_n \le 1 - \frac{\tau}{n}\right) \to e^{-\tau} $$
$$ P\left(n(1 - M_n) < \tau\right) = 1 - e^{-\tau} + o(1) $$
\end{example}

\begin{example}[Индикатор событий]
Пусть $X_i = \begin{cases} 1, & \xi_i > v_n \\ 0, & \xi_i \le v_n \end{cases}$, где $P(X_i = 1) = 1 - F(v_n)$.

Число превышений $S_n = \sum X_i$ имеет биномиальное распределение:
$$ P(S_n = k) = \binom{n}{k}(1 - F(v_n))^k (F(v_n))^{n-k} $$

Задача сводится к распределению Бернулли.
\end{example}
